\section{Discussion and limitations}
\label{sec:disc}

\paragraph{What the evidence supports.}
(1)~A handful of well-chosen benchmarks gives most of what a larger budget would: on MTEB-main five tasks of $\sim$215 explain 71\% of the variance (relative to the mean baseline) of a new model's other scores, ten explain 76\%, and no budget up to 40 passes $\approx$80\%. (2)~The choice of benchmarks matters mainly below $k\approx10$ and the gain over random is substantial there (13--21\% in MSE), but it is not uniformly positive: it is absent for homogeneous pools and for $k$ comparable to the pool size, and it is negative for LiveBench-7 at $k=3$. (3)~Heuristics that look reasonable can hurt: the most-evaluated benchmarks are worse than random at every $k$ we tried on MTEB-main, and the most-correlated ones from $k=6$ on. (4)~The theory explains the shape: an irreducible per-benchmark floor, a requirement $k\gtrsim r(1+1/\epsilon)$ for a factor structure of rank $r$, and a spectral lower bound that no selection can beat; the sandwich and the finite-history inflation are accurate in simulation, and the fitted Gaussian predicts its own error on real matrices. The hypothesis that a small subset suffices to infer the rest is thus \emph{supported in the regime of moderate accuracy} (about 75\% of variance) and \emph{not supported for accurate prediction of individual scores}, where a plateau remains.

\paragraph{Where it fails or is unproven.}
\textbf{Plateau.} A residual of 20--30\% of variance is benchmark-specific; we do not know how much is evaluation noise and how much true specialization, so we cannot say whether more data or better predictors would remove it. \textbf{Gaussian linear model.} The mean-and-covariance model ignores ceiling effects and non-linear dependence; \textsc{kNN} beats it at $k=1$ by 5--8 points of $R^2$, and a logit transform gave no gain. The factor-model covariance, although natural for the theory, selected 5--9\% worse than the pairwise covariance at $k\ge10$, and low-rank prediction was never better than the full covariance: our theory is a guide to scaling, not a description of these matrices. \textbf{Selection is not adaptive.} Under the Gaussian model adaptivity gives nothing (Prop.~\ref{prop:closed}), but we did not test adaptive policies with non-Gaussian predictors. \textbf{Retrospective protocol.} A model's observed benchmarks are not random: authors choose where to report (missing not at random); we restrict selection to the model's observed set and score on its other observed benchmarks, which may flatter methods in either direction. Costs per benchmark are ignored; a cost-aware greedy ratio rule is a one-line change but untested. \textbf{Data.} MTEB models are text-embedding models, not generative LLMs; the three LLM matrices have $N\le22$, so $k$ cannot exceed 3--12 and conclusions about large budgets rest on MTEB alone. Near-duplicate models (merges, fine-tunes, sizes of one family) make random splits optimistic; the organisation and temporal splits reduce but do not remove this. The time split for MTEB has only 53 test models. The MTEB subset we use consists of tasks evaluated for $\ge$100 models, which over-represents widely run tasks. We tuned hyper-parameters on the training rows only, but chose \textsc{Gauss} as the default after looking at inner validation on MTEB-main, and the study designed its protocol after a first, flawed run with complement targets (\S\ref{sec:exp_pitfall}; that section reports a controlled re-run, not the first run). \textbf{Scope of the claims.} Statements about ``$k$ needed'' depend on the target accuracy, the matrix, and the shift between history and new model; they are not universal constants.

\paragraph{Practical guidance.} With a history like MTEB-main, run 5--10 benchmarks chosen by greedy variance reduction (or $k$-medoids), predict with the conditional Gaussian (\textsc{kNN} for $k=1$), expect the \emph{average} of the other scores to be within $\sim$1--2 points but individual scores only within $\sim$7--8, plan $k$ with the model-implied error, and re-check against full evaluations when the new model comes from a different lab or period.

\paragraph{Broader impact.} Cheaper evaluation saves energy and lowers entry barriers, but risks over-reliance on proxy benchmarks; inferred scores should be labelled as such.
